% AI-assisted manuscript: ChatGPT/Codex, OpenAI, 2026-09-23.
% Exact model version/release date must be confirmed from the team's interface.
\documentclass[UTF8,zihao=-4,fontset=fandol,a4paper]{ctexart}
\usepackage[left=25mm,right=25mm,top=23mm,bottom=23mm]{geometry}
\usepackage{amsmath,amssymb,mathtools}
\usepackage{unicode-math}
\setmathfont{Latin Modern Math}
\usepackage{graphicx,booktabs,longtable,tabularx,array}
\usepackage{caption,float,placeins,enumitem}
\usepackage[hidelinks]{hyperref}
\usepackage{xurl}
\setmainfont{Latin Modern Roman}
\setlength{\parindent}{2em}
\setlength{\parskip}{0pt}
\linespread{1}
\ctexset{section={format=\centering\heiti\zihao{4},beforeskip=14pt,afterskip=8pt},subsection={format=\songti\bfseries\zihao{-4},beforeskip=9pt,afterskip=5pt},subsubsection={format=\songti\bfseries\zihao{-4}}}
\captionsetup{font=normalsize,labelfont=bf,labelsep=quad,skip=6pt}
\setlist{nosep,leftmargin=2em}
\renewcommand{\arraystretch}{1.15}
\setlength{\tabcolsep}{5pt}
\setlength{\textfloatsep}{12pt plus 2pt minus 2pt}
\setlength{\floatsep}{10pt plus 2pt minus 2pt}
\setlength{\intextsep}{10pt plus 2pt minus 2pt}
\renewcommand{\topfraction}{.88}
\renewcommand{\bottomfraction}{.75}
\renewcommand{\textfraction}{.08}
\renewcommand{\floatpagefraction}{.75}
\pagestyle{plain}
\hypersetup{pdftitle={地形与能量约束下的山区无人机运输、通信与资源分区},pdfauthor={},pdfsubject={2026年中国研究生数学建模竞赛D题匿名正文}}
% Auto-generated; never edit numbers here.
\newcommand{\MassTightBatches}{5}
\newcommand{\VolumeTightBatches}{9}
\newcommand{\EnergyTightBatches}{4}
\newcommand{\NumBoxes}{80}
\newcommand{\NumServices}{15}
\newcommand{\QOneTrips}{18}
\newcommand{\QOneEnergy}{59.130}
\newcommand{\QOneTime}{32776.016}
\newcommand{\QOneHours}{9.104}
\newcommand{\QTwoTrips}{24}
\newcommand{\QTwoEnergy}{65.566}
\newcommand{\QTwoTime}{5903.038}
\newcommand{\QTwoHours}{1.640}
\newcommand{\QTwoLate}{0.000}
\newcommand{\QTwoFixedBound}{98.149}
\newcommand{\QTwoOrderGain}{14.10}
\newcommand{\QThreeTrips}{24}
\newcommand{\RelayTrips}{3}
\newcommand{\QThreeTransportEnergy}{66.482}
\newcommand{\RelayEnergy}{2.659}
\newcommand{\JointEnergy}{69.141}
\newcommand{\JointTime}{5931.735}
\newcommand{\JointHours}{1.648}
\newcommand{\QThreeTransportTime}{5931.735}
\newcommand{\QThreeLate}{0.000}
\newcommand{\HardSlack}{266.282}
\newcommand{\HardBoxes}{31}
\newcommand{\StatesOne}{35061}
\newcommand{\StatesHalf}{69843}
\newcommand{\MarginOne}{0.091}
\newcommand{\MarginHalf}{0.091}
\newcommand{\StatesQuarter}{139402}
\newcommand{\MarginQuarter}{0.091}
\newcommand{\ViolationsQuarter}{0}
\newcommand{\RepairLateReduction}{100.00}
\newcommand{\RepairEnergyReduction}{26.61}
\newcommand{\RepairTimeIncrease}{-4793.12}
\newcommand{\ViolationsHalf}{0}
\newcommand{\CriticalTripHalf}{column-019}
\newcommand{\CriticalTimeHalf}{4684.446}
\newcommand{\ParetoCount}{42}
\newcommand{\SearchSeed}{20260925}
\newcommand{\QTwoIterations}{110000}
\newcommand{\QThreeIterations}{5000}
\newcommand{\ReserveTripBreak}{23.09}
\newcommand{\ReserveReachabilityLimit}{35.27}
\newcommand{\MinSoc}{20.10}
\newcommand{\QOneReserveLimit}{23.09}
\newcommand{\QThreeEnergyTolerance}{0.12}
\newcommand{\QTwoEnergyTolerance}{4.02}
\newcommand{\GapTwo}{1}
\newcommand{\GapThree}{3}
\newcommand{\EnergyIncrease}{5.45}
\newcommand{\TimeIncrease}{0.49}
\newcommand{\CMinPayload}{58.91}
\newcommand{\BMinPayload}{28.80}
\newcommand{\CRestricted}{5}
\newcommand{\CandidateCount}{9}
\newcommand{\CandidateMaxTrips}{25}
\newcommand{\CandidateMinTrips}{17}
\newcommand{\AtomicGroups}{9}
\newcommand{\PartitionsTwo}{255}
\newcommand{\PartitionsThree}{3025}
\newcommand{\ReserveGridCount}{61}
\newcommand{\CapacityCaseCount}{2745}
\newcommand{\BatchScenarioCount}{5}
\newcommand{\MaxRouteServices}{5}
\newcommand{\MergeLateScale}{200000}
\newcommand{\MergeTimeScale}{10000}
\newcommand{\MergeEnergyScale}{100}
\newcommand{\RelayGuard}{5}

\newcommand{\paperfig}[3]{\par\begin{figure}[htbp]\centering\includegraphics[width=\linewidth]{../outputs/figures/#1.pdf}\caption{#2}\label{fig:#3}\end{figure}\par}
\newcommand{\papertable}[3]{\par\begin{table}[htbp]\centering\caption{#2}\label{tab:#3}\input{tables/#1}\end{table}\par}
\newcommand{\longresulttable}[3]{\begingroup\def\resultcaption{#2}\def\resultlabel{tab:#3}\input{tables/#1_long}\endgroup}
\begin{document}
\pagenumbering{arabic}
\begin{center}
{\heiti\zihao{3}地形与能量约束下的山区无人机运输、通信与资源分区}\par
\vspace{12pt}
{\heiti\zihao{4}摘要}\par
\end{center}
针对山区 \NumServices{} 个服务区、\NumBoxes{} 个不可拆货箱的运输任务，本文以 DEM 航高、动态剩余载荷和分阶段能量构成统一物理底座，将能力边界、资源调度、通信时序和独立配置连成递进模型。

针对 Q1，先扣除地形爬升、空载返程与返航储备，利用载荷—能耗单调性二分求安全载荷，再对满足质量、体积和能量约束的货箱子集进行集合划分动态规划。得到 \QOneTrips{} 架次、\QOneEnergy{} kWh、累计作业 \QOneTime{} s；精确性限于单点往返及“架次—能耗—累计时间”词典序目标。

针对 Q2，加入实体无人机、共享电池充电、硬截止与软延误，采用紧迫任务优先构造、可行合并和资源重排，保存 \CandidateCount{} 个有限候选。显式选择 \QTwoTrips{} 架次方案，能耗 \QTwoEnergy{} kWh，完工 \QTwoTime{} s，全部硬截止满足；相邻候选比较说明继续压缩架次须付出及时性代价，该选择不构成全局最优证明。

针对 Q3，按双向预算与 DEM 遮挡检查飞行、交接全过程，以指定单跳中继覆盖直连不足区间。在固定分组、路线、资源顺序和中继位置/高度下，用词典序 LP 优化运输开始和中继窗口，得到 \QThreeTrips{} 个运输架次、\RelayTrips{} 个中继架次，联合能耗 \JointEnergy{} kWh、完工 \JointTime{} s。1 s 与 0.5 s 回放均未发现中断，最小硬截止余量 \HardSlack{} s；最优性仅针对固定方案的时序问题。

针对 Q4，继承 Q3 的任务与时序，将服务区归并为 \AtomicGroups{} 个不可拆连通分量，以区间峰值核算资源，精确枚举两组和三组分区，缺口分别为 \GapTwo{} 与 \GapThree{}。该结果以组内同型资源可重编号、中继完整任务可按组复制为前提。

输入摘要核验、逐箱逐资源回放与边界反例构成验证链。但加密通信最小余量仅 \MarginHalf{} dB，固定运输计划的返航电量仅容许能耗统一上浮约 \QThreeEnergyTolerance\%。结果支持官方确定性场景中的可复现候选，不能替代连续通信证明或实飞验证。

\vspace{8pt}
\noindent\textbf{关键词：}山区无人机运输；动态载荷能耗；集合划分；共享电池调度；通信中继；独立资源配置

\clearpage
\section{问题分析与统一建模框架}
\subsection{从运载能力到可执行方案}
山区应急投送的关键困难，是一条地理上可达的航线未必能在给定电量、交付时限和通信条件下执行。题目提供一个调度中心、\NumServices{} 个服务区与 \NumBoxes{} 个不可拆分货箱，要求从单点往返逐步扩展至异构运输调度、空中中继和独立任务分区\cite{problem}。本文以官方标准测试场景为研究对象，不把背景灾情叙述直接作为需求预测数据。

运输路径与能量消耗不能分开处理。既有无人机配送研究已指出，载荷相关能耗和多架次复用会影响路线可行性\cite{dorling}；不同能耗模型的使用条件也不能互相替代\cite{zhang}。本题的重点不是重新拟合动力学，而是在官方统一规则下，让每个阶段的载荷、高度、用电、充电和通信判定保持一致。

我们把四问组织为一条逐步增加约束的计算链。Q1 给出不受实体资源调度影响的能力基线；Q2 加入无人机、电池和交付时限，得到可执行的运输候选；Q3 在固定离散候选上联动运输开始时刻与中继服务窗口；Q4 保持 Q3 任务不变，度量跨组资源共享被取消的代价。四问的可行域不同，因此架次数和能耗的差异不能直接解释为同一优化问题的算法优劣。

\subsection{统一思想与证据层级}
统一框架如图\ref{fig:framework}。底层把每个航段转化为可复用的时间、能量和轨迹对象，上层只决定货箱、路径、资源与时序。相同的物理对象再交给验证器重算，避免求解模块自己报告成功后直接进入论文。
\paperfig{fig01_framework}{同一物理计算连接能力边界与资源、通信时序。左侧为真实 DEM 剖面和安全载荷曲线，右侧为实际运输与中继资源占用；中继浅色条表示飞行与准备，深色条表示通信服务。Q1 为精确组批，Q2 为启发式候选，Q3 的线性规划最优性仅针对固定离散方案。}{framework}

\FloatBarrier

\section{假设、数据处理与符号}
\subsection{题设口径与补充假设}
所有需求、时限和设备参数取自官方五份工作簿，地形取官方 30 m DEM。坐标为经纬度，水平距离采用 WGS84 椭球测地线距离\cite{karney}；航迹与 DEM 穿越按代码一致的经纬度端点线性插值计算。该短距离近似保持各模块口径一致，但不代表高精度航空导航轨迹。

依据题设，货箱不可拆分，每箱恰好交付一次；每个运输架次由 O01 出发并最终返回 O01。交付后对应货箱质量从剩余载荷中扣除。下降不单计附加能量，运输交接不另加未经题设给定的悬停功率。中继建链阶段按悬停与通信功率计入能耗。所有能源资源初始满电，再次投入前充至满电，不同资源可并行充电。

赛题未提供风速、风向或阵风参数，故采用给定计划速度，不另建风修正模型。DEM 视为像元内分段常数；线路恰好经过边界或角点时，纳入全部相接像元，以较高地形进行保守检查。中继窗口在已发现的直连不足区间两侧各扩展 5 s，这是明确的时序保护假设，不是连续通信证明。

\subsection{数据检查与符号系统}
读取后先检查编号唯一性、货箱归属、资源类型和数值范围，再检查必需航段和通信射线是否超出 DEM 或经过无效像元。正式计算使用完整精度；论文的三位能耗小数、两位载荷小数和秒级时间展示仅用于阅读。关键运输参数见表\ref{tab:parameters}，完整参数保留于原始工作簿。
\papertable{drone_parameters}{三种运输机型的主要参数；电池可用能量与额定载荷分别约束航程和装载}{parameters}

集合 $\mathcal S$ 表示服务区，$\mathcal B$ 表示货箱，$\mathcal G$ 表示机型，$\mathcal P$ 和 $\mathcal R$ 分别表示运输及中继架次。$i,j$ 为节点，$b$ 为货箱，$g$ 为机型，$p,r$ 为架次索引。表\ref{tab:symbols}区分参数、决策与派生状态。
\begin{table}[htbp]\centering\caption{核心符号及单位}\label{tab:symbols}
\begin{tabularx}{\linewidth}{lXl}\toprule
符号 & 含义与类别 & 单位\\\midrule
$d_{ij},z_i,H_{ij}$ & 水平距离、节点作业海拔、计划巡航海拔 & m\\
$m_b,v_b,w_b$ & 货箱质量、体积、应急优先系数（参数） & kg，m$^3$，1\\
$Q_g,V_g,E_g,\rho_g$ & 额定载荷、体积、电池可用能量、返航储备 & kg，m$^3$，kWh，1\\
$q_{p,ij},s_p$ & 航段剩余载荷（状态）、架次开始时刻（决策） & kg，s\\
$x_{pb},u_{pd},v_{pk}$ & 货箱、实体无人机、电池的分配变量 & 0--1\\
$C_b,\ell_b,T$ & 交付完成时刻、延误、任务完工历时（派生量） & s\\
$a_r,b_r,\bm h_r$ & 中继建链完成、服务结束、三维悬停位置 & s，s，经纬度/m\\
$M_{ij}(t)$ & 双向链路预算减去实际传播损耗（余量） & dB\\\bottomrule
\end{tabularx}\end{table}
\FloatBarrier

\section{地形—载荷—能量统一计算模型}\label{sec:physics}
四问要在同一标准场景下衔接，首先必须让同一条路线得到相同的时间、能耗和状态序列。本文将运输架次表示为“货箱集合、机型、访问顺序”，公共引擎据此返回逐段载荷、飞行阶段、交付偏移、总时长与总能耗。上层算法决定怎样组批和排程，不能改变这套物理计算。

\subsection{DEM 航高与飞行阶段}
山区航线的代价并不只由水平距离决定：沿途一处高地就会抬升整个航段的计划巡航海拔。记线路经过的闭 DEM 像元集合为 $\mathcal C_{ij}$，像元高程为 $h_c$；O01 作业海拔取地面高程，服务区作业海拔取地面以上 30 m。对不同节点间航段，有
\begin{align}
H_{ij}&=\max_{c\in\mathcal C_{ij}}h_c+50,\qquad
h^+_{ij}=H_{ij}-z_i,\quad h^-_{ij}=H_{ij}-z_j,\label{eq:altitude}\\
\tau_{gij}&=\frac{h^+_{ij}}{v_g^\uparrow}+\frac{d_{ij}}{v_g^c}+\frac{h^-_{ij}}{v_g^\downarrow}.\label{eq:flight}
\end{align}
长度以 m、速度以 m/s 计，所得时间为 s。每次完成服务区交付后，下一航段从该点作业高度重新爬升；不能把多点配送当成一次起降、全程等高巡航。同一节点到自身的零航段则取零距离、零升降。
\paperfig{fig03_physics}{同一水平航段的运输代价仍随地形航高和载荷改变，空载返程也消耗能源。S014 的真实 DEM 剖面对应上图，实际 Q1 批次的去程载货与返程空载能量分解对应下图；下降附加能耗为零遵循题设规则。}{physics}
图\ref{fig:physics}把航高与能量联系起来：终点较低并不意味着沿途无需高爬升，空载返程也不能省略。因而相同水平距离可对应不同的山区运输代价，后续载荷边界须逐航段计算。

\subsection{动态载荷与等效航程}
载荷相关能耗会改变配送路线的可行域\cite{dorling}，但不同能耗模型具有各自适用条件\cite{zhang}。本题已有统一计算规则，故直接使用官方等效航程，不再用外部经验模型替换。机型 $g$ 携带有效载荷 $q$ 时，
\begin{equation}
L_g(q)=L_g^0-(L_g^0-L_g^F)(q/Q_g)^{3/2},\qquad 0\le q\le Q_g.\label{eq:range}
\end{equation}
对架次 $p$，令 $\mathcal B_p$ 为其货箱集合，$\pi_p$ 为从 O01 出发并返回的节点序列。进入航段 $i\to j$ 时，
\begin{equation}
q_{p,ij}=\sum_{b\in\mathcal B_p:\,b\text{尚未完成交付}}m_b.\label{eq:payloadstate}
\end{equation}
每个服务区完成交接后移除对应货箱，再计算下一段。由此，载荷不是“起飞质量一直保持不变”，也不是“每段都按空载”；全部货物交付后的返航段才取 $q=0$。

\subsection{分阶段能量与返航储备}
水平航程和正向爬升分别消耗电量：
\begin{equation}
E_{gij}(q)=E_g\frac{d_{ij}}{L_g(q)}+
\frac{(m_g^0+q)g_0h^+_{ij}}{\eta_g\,3.6\times10^6}.\label{eq:energy}
\end{equation}
其中 $m_g^0$ 为含电池空载总质量，$g_0=9.80665$ m/s$^2$，$\eta_g$ 为爬升效率。第二项由焦耳换算为 kWh，与第一项一致。下降附加项按题设取零，并不意味着真实下降过程不耗电，而是本题等效模型的规定。

架次可行性同时要求
\begin{equation}
\begin{gathered}
\sum_{b\in\mathcal B_p}m_b\le Q_{g_p},\qquad
\sum_{b\in\mathcal B_p}v_b\le V_{g_p},\\
E_p=\sum_{(i,j)\in\pi_p}E_{g_pij}(q_{p,ij})
\le(1-\rho_{g_p})E_{g_p}.\label{eq:feasible}
\end{gathered}
\end{equation}
质量限制回答“是否超载”，体积限制回答“能否装入”，能源限制回答“能否完成全程并留出返航储备”。由于只卸货、不取货，起飞质量和体积是全程最大值；电量则必须累计所有航段。共享这一判据，比单纯要求四问使用相同参数名称更重要。

\subsection{无人机与能源资源的独立占用}
作业从固定准备和逐箱装载开始。各段飞行之后，按服务区基础交接时间及逐箱附加时间计算完成时刻；货箱的交付偏移 $\delta_{pb}$ 取交接结束，不能用到达代替。公共引擎最终给出 $(t_p,E_p,\delta_{pb})$，Q2、Q3 再通过 $C_b=s_p+\delta_{pb}$ 转成绝对时刻。

无人机返航即可结束该运输架次占用；电池还须充至满电。令 $\sigma_p=1-E_p/E_{g_p}$，官方两阶段充电时间为
\begin{equation}
 t_{\rm chg}(\sigma)=T_{\rm full}
 \begin{cases}
 0.65(0.90-\sigma)/0.90+0.35,&0\le\sigma<0.90,\\
 0.35(1-\sigma)/0.10,&0.90\le\sigma\le1.
 \end{cases}\label{eq:charge}
\end{equation}
于是无人机和电池占用区间分别为 $[s_p,s_p+t_p)$ 与 $[s_p,s_p+t_p+t_{\rm chg}(\sigma_p))$。半开区间允许资源在前一任务释放的时刻立即投入下一任务。它们将在 Q2 约束排程，在 Q3 约束时序 LP，并在 Q4 转为独立资源需求，构成四问之间可核验的公共接口。
\FloatBarrier

\section{Q1：最大安全载荷与精确组批}\label{sec:q1}
\subsection{为什么额定载荷不足以判断安全}
额定载荷只限制装载质量，不保证山区往返能源充足。载货去程既要完成水平飞行，又要把整机与货物抬升至 DEM 决定的巡航高度；卸货后的返程仍有飞行和爬升开销。还须保留题设返航电量，因此应先回答每个服务区—机型组合“最多可安全运多少”，再讨论怎样把实际货箱装成批次。

\subsubsection{服务区—机型安全载荷边界}
Q1 限定单点往返 $O01\to i\to O01$。定义
\begin{equation}
q^*_{gi}=\max\{q\in[0,Q_g]:
E_{g,0i}(q)+E_{g,i0}(0)\le(1-\rho_g)E_g\}.\label{eq:qmax}
\end{equation}
该边界同时扣除了空载返程和储备，不能仅用额定载荷或单程航程判断。若空载往返已不可行，则该组合记为不可达；若额定载荷可行，则 $q^*_{gi}=Q_g$；其余情况求能量边界。不可达与“可达但只能承载零质量”是不同状态，输出中分别标记。

\subsubsection{单调性与二分依据}
当官方参数满足 $L_g^0>L_g^F>0$ 时，式\eqref{eq:range}对 $q$ 非增，故水平项 $E_gd/L_g(q)$ 非减，爬升项也随 $q$ 非减；空载返程为常量。因此载荷可行集是一个区间，二分能定位其上界。提高 $\rho_g$ 只会缩小能量预算，不能提高安全载荷，这是数值求解与边界测试共用的单调性依据。
\paperfig{fig03_payload_boundary}{额定载荷与能量共同形成服务区—机型能力边界。上图列出全部组合，下图给出安全载荷占额定值比例；C 型部分远端服务区先触及能源限制。各点对应不同地形，连线不能解释为控制其他因素后的纯距离响应。}{payload}
结果如图\ref{fig:payload}：C 型有 \CRestricted{} 个服务区不能满额定载荷往返，其最小安全载荷为 \CMinPayload{} kg；B 型最小为 \BMinPayload{} kg。A 型在基准储备下达到额定载荷。这表明机型“大”不等于任意服务区都能满载，能力判断必须同时考虑其能量与航程参数。

\subsection{不可拆货箱的集合划分模型}
安全载荷是连续质量边界，而需求由带有编号、质量和体积的不可拆货箱构成。即使总质量可运，也未必能按平均质量切成可行批次。对服务区 $i$ 的货箱集合 $\mathcal B_i$，枚举非空子集 $S$ 和机型 $g$，仅保留同时满足式\eqref{eq:feasible}的单点往返候选。质量和能源相关却不等价：前者限制起飞载荷，后者还受 DEM 和返程影响；体积更不能由质量替代。

Q1 不要求实体无人机分配或交付时限，故先最少架次，在同架次数内再减少能耗，最后减少累计作业时间：
\begin{equation}
\operatorname{lexmin}\left(N_1,\sum_pE_p,\sum_pt_p\right).\label{eq:q1objective}
\end{equation}
这一优先关系既响应单点运输能力问题，也避免以任意权重把“少一次飞行”与“少一点电量”相抵消。累计时间包含准备、装载、往返和交接，是作业量指标，尚未表示多机并行完工。

\subsection{锚定子集动态规划与精确性}
本题每个服务区的货箱规模允许枚举子集，且不同服务区在 Q1 中不共享调度约束，因而适合用集合划分动态规划，不需要引入更复杂的全局调度搜索。令 $D_i(U)$ 表示已交付集合 $U\subseteq\mathcal B_i$ 的最佳目标三元组，取 $D_i(\varnothing)=(0,0,0)$，其余状态初始不可达。对可行且与 $U$ 不交的候选 $S$，
\begin{equation}
D_i(U\cup S)\leftarrow\min_{\rm lex}
\{D_i(U\cup S),\ D_i(U)+(1,E_g(S),t_g(S))\}.\label{eq:dp}
\end{equation}
转移表示新增一架次恰好交付 $S$，终态 $D_i(\mathcal B_i)$ 即覆盖全部货箱且无重复的最优划分。相同 $S$ 只保留能耗最小、同能耗时时间最短的机型：它们覆盖的箱子相同，又无实体资源耦合，故删去被支配候选不影响 Q1 目标。

递推只选含最小编号未交付货箱的下一批次。每个完整划分必有一批包含该货箱，按该批先处理只是消除批次排列重复，不遗漏划分。各区分别精确求解后将目标相加，得到本问题定义下的词典序精确解。若一个服务区有 $n_i$ 箱，状态不超过 $2^{n_i}$，最坏转移规模为 $O(3^{n_i})$；精确性不意味着它能直接解决大规模实体调度。

\subsection{组批结果、约束余量与小结}
\textbf{Q1 得到 \QOneTrips{} 架次、\QOneEnergy{} kWh 和累计作业 \QOneTime{} s。}全部 \NumBoxes{} 箱恰好交付一次，最大安全载荷与完整批次见附录。图\ref{fig:batching}将箱级分配与约束占用放在一起，说明每一批为什么可行，而非只展示架次数。
\paperfig{fig04_batching}{不可拆货箱通过精确划分形成 \QOneTrips{} 个可行批次。条宽为质量、条内数字为箱数、纹理区分货类；右侧核对质量、体积和能量预算。这里展示批次，不是尚未定义的实体无人机时间表。}{batching}
验证逐批重新累加质量与体积、计算载货去程和空载返程。固定组批最低返航 SOC 为 \QOneReserveLimit\%，高于基准储备；储备改变后的重新组批见第\ref{sec:sensitivity}节。精确最优的成立范围是单点往返、给定物理规则、当前数值容差及式\eqref{eq:q1objective}，不含后续新增的实体资源、时限和通信条件。

\textbf{本问小结。}我们得到服务区—机型安全载荷边界与不可拆货箱的精确组批基准；单调性、完整子集划分和逐箱回放使其可信。但它尚不能回答哪架飞机、哪块电池何时执行。Q2 继承统一能量判据和能力认识，重新选择多点分组与排程，既不强制沿用 Q1 批次，也不能直接用单点往返载荷界代替多点路线的逐段能源检查。
\FloatBarrier

\section{异构运输调度与候选方案权衡}
\subsection{时限与资源模型}
Q2 在同一物理模型上增加实体无人机、电池及配送时限。定义 $x_{pb}=1$ 表示货箱 $b$ 由架次 $p$ 承运，且 $\sum_px_{pb}=1$。机型、访问顺序确定后，公共航段引擎给出作业时长 $t_p$、能耗 $E_p$ 和货箱交付偏移 $\delta_{pb}$，因而
\begin{equation}
C_b=s_p+\delta_{pb},\quad
\ell_b\ge C_b-e_b,\quad\ell_b\ge0,\quad
C_b\le d_b\quad(b\in\mathcal B_{\rm hard}).
\label{eq:deadline}
\end{equation}
$e_b$ 是期望时间，$d_b$ 是硬截止时间：医疗物资以期望时间作为硬约束，首批保障以首批截止为硬约束，同时满足两种身份时取更早者。其他物资期望时间进入加权延误 $W=\sum_bw_b\ell_b$。

若架次 $p,p'$ 共用无人机，则区间 $[s_p,s_p+t_p)$ 不得重叠；若共用电池，则区间延伸至 $s_p+t_p+t_{\rm chg}(1-E_p/E_g)$。把电池与无人机绑定成不可拆的固定组合会丢失共享电池的调度能力，因此分别记录两类分配变量。多指标为
\begin{equation}
\min\left(W,\ T_2,\ E_2,\ N_2\right),\qquad
T_2=\max_p(s_p+t_p),\quad E_2=\sum_pE_p.
\label{eq:q2objective}
\end{equation}
式\eqref{eq:q2objective}表达权衡关系；实现中通过有限候选搜索和人工明确选点处理，并未求出完整 Pareto 前沿。

\subsection{确定性启发式与选点原则}
首先以紧迫货箱为锚构造满足硬时限的路线，优先保障硬时限，再综合延误、承运质量和箱数。其次尝试合并架次，枚举候选访问顺序并重新选择可用资源、重算出发时间；每条路线最多枚举五个服务区，以限制阶乘复杂度。每次合并后的载荷、逐箱完成时间、能源和资源占用都重新计算。

保存架次数从 24 至 17 的八个候选，而不是只留下选中的结果。附录表\ref{tab:frontier}与图\ref{fig:optimization}表明，22 架次相较更少架次保留了较短完工历时和较低延误；相比 23 架次，它降低能耗，但增加部分软时限延误。因此选用 \QTwoTrips{} 架次是一个显式折中，不能称为由所有目标共同唯一决定的最优点。
\paperfig{fig05_optimization}{较少架次需要付出及时性代价。上图展示八个实际候选的能耗与加权延误，折线仅为这两个显示目标下的非支配关系，星号标记所选方案；下图按实体无人机展示该方案的占用区间。有限候选集不构成全局前沿证明。}{optimization}

\subsection{结果及其含义}
\textbf{所选方案用 \QTwoTrips{} 架次完成全部货箱，运输能耗 \QTwoEnergy{} kWh，完工历时 \QTwoTime{} s。}加权延误为 \QTwoLate{} 优先系数·s，全部硬截止要求满足。逐箱交付表将到达与完成交接严格区分，资源检查也覆盖电池返航后的充电阶段。

图中的实体无人机占用并不均衡：大载荷机型承担较多后续任务，小型机较早释放。这是当前候选下的任务分配结果，不应把长空闲直接归因于算法错误；但它提示还可扩大跨机型交换和离散路线搜索，以评估剩余改进空间。本文没有该问题的全局下界，故不报告未经证实的最优性差距。
\FloatBarrier

\section{通信保障下的运输与中继联合时序}
\subsection{双向链路与单跳中继}
空中中继可以补充灾后地面覆盖，但其位置、回传链路与服务时段必须同时可行\cite{zeng}。按官方通信规则，双向允许传播损耗取两个方向预算的较小值。对三维端点间距离 $D_{ij}(t)$（km）和频率 $f$（MHz），定义
\begin{align}
L_{ij}(t)&=32.45+20\log_{10}f+20\log_{10}D_{ij}(t)+L_{\rm obs}\,b_{ij}(t),\label{eq:fspl}\\
M_{ij}(t)&=L^{\max}_{i\leftrightarrow j}-L_{ij}(t).\label{eq:margin}
\end{align}
自由空间项采用与 ITU-R P.525 一致的频率、距离量纲\cite{itu}。$b_{ij}=1$ 表示视线射线与闭 DEM 像元地形相交，包括切触；有遮挡时附加题设损耗。官方参数给出 2400 MHz、遮挡附加 10 dB，直连、接入和回传预算分别为 122、116、126 dB。

令 $\mathcal R(t)$ 为已完成建链且尚未结束服务的中继集合，则运输状态可用的必要判据写为
\begin{equation}
M_p(t)=\max\left\{M_{pG}(t),\ \max_{r\in\mathcal R(t)}\min[M_{pr}(t),M_{rG}(t)]\right\}\ge0.
\label{eq:comm}
\end{equation}
实现中逐运输架次使用其配置的单个保障中继，不能把不可同时工作的不同中继的链路拼接，也不允许中继之间多跳。直连足够时可直接连接网关；需要中继时接入、回传与时间窗口必须同时成立。

\subsection{中继能量与固定离散方案的线性规划}
一个中继架次的能耗由往返巡航、正向爬升、建链与服务组成：
\begin{equation}
E_r=E_r^{\rm cruise}+E_r^{\rm climb}
 +(P_r^{\rm hover}+P_r^{\rm comm})\frac{t_r^{\rm setup}+b_r-a_r}{3600},
\qquad E_r\le(1-\rho_r)E_r^{\rm use}.
\label{eq:relay}
\end{equation}
其中两个功率以 kW 计，服务与建链时间以 s 计。中继悬停点位于 DEM 内且离地不超过 300 m。来回航段既满足地形净空，也到达选定悬停海拔；建链完成前不计为可提供服务。

配置文件固定货箱、路径、机型、资源先后关系以及中继坐标。对该离散方案，逐相位采样得到需要中继的相对时间范围 $[\alpha_p,\beta_p]$，加入
\begin{equation}
a_{r(p)}\le s_p+\alpha_p-5,\qquad
b_{r(p)}\ge s_p+\beta_p+5.
\label{eq:window}
\end{equation}
再加入式\eqref{eq:deadline}、资源顺序、建链前飞行时间、中继周转时间和式\eqref{eq:relay}。由于航段及载荷固定，上述约束对开始时刻和服务窗口均为线性。依次最小化加权延误、联合完工历时、中继服务总时长，前一层最优值以数值容差锁定后求下一层，使用 HiGHS 线性规划求解器\cite{highs}。运输能耗此时固定，中继服务时长最小化也就减少了可调整的联合能耗。

这一步只证明给定离散方案下的时序最优性。悬停点的全空间搜索、货箱组批和资源顺序的联合全局最优性没有由该 LP 证明。论文保留这一条件，防止把“LP 已最优”扩大解释为 Q3 原问题全局最优。

\subsection{联合方案与通信保障代价}
\textbf{联合候选含 \QThreeTrips{} 个运输架次与 \RelayTrips{} 个中继架次，总能耗 \JointEnergy{} kWh，联合完工 \JointTime{} s。}其中运输能耗 \QThreeTransportEnergy{} kWh，中继能耗 \RelayEnergy{} kWh；运输完工时刻为 \QThreeTransportTime{} s。加权延误为 \QThreeLate{} 优先系数·s，硬截止最小余量 \HardSlack{} s。
\paperfig{fig06_communication}{中继覆盖实际轨迹上直连不足的采样状态。橙点表示直连不足，空心三角为中继悬停点，星号为最紧通信状态；下图为 60 s 时间窗内全部在途状态的最小可用链路余量。检查包含每个飞行与交接相位的端点，共 \StatesOne{} 个状态，最小余量 \MarginOne{} dB；该结论仅针对采样集合。}{communication}

与选定 Q2 候选相比，Q3 联合能耗增加 \EnergyIncrease\%，完工历时增加 \TimeIncrease\%。这是当前两个候选在不同约束集合下的结果差异，不是通信约束的理论最小代价。图\ref{fig:communication}显示，通信困难并非只发生在目的节点；只检查起点、终点和中继位置会漏掉航段中的遮挡状态。

射线遮挡是离散地形上的非光滑判定，通信余量很小的区间尤其需要谨慎。尽管所有已检查状态通过，仍不能据此断言“任意连续时刻均满足题设”。第\ref{sec:validation}节进一步加密采样并报告余量变化。
\FloatBarrier

\section{Q4：独立任务分区与资源配置}\label{sec:q4}
Q3 的正式候选依赖全体任务共享无人机与能源资源。本问继承通信修复后重新计算并通过采样验证的正式输出，逐项核对任务、保障关系和时间区间；配置结论仍以这些固定安排为条件。若救援被划为独立执行组，即使货物、路线和时刻都不变，原本能相互周转的资源也可能不能跨组调用。Q4 因而研究固定计划的独立执行代价，不重新寻找运输路线或通信方案。

\subsection{从 Q3 继承任务与不可拆连通分量}
将同一运输架次访问的服务区连接；由传递闭包得到 \AtomicGroups{} 个不可拆服务区连通分量。一个架次不能分给两个组，所以任务组只能由这些分量合并形成，每个分量恰归一组且各组非空。此时可分的基本单元已不是单个服务区，更不是单个货箱。

同时保持 Q3 的运输时间与中继保障关系。若一个中继架次原来服务的运输任务被分到不同组，则在各组复制该中继架次的完整位置和服务窗口，并分别计入中继无人机、能源组件及能耗。同一组内多条运输任务共用同一中继架次时仅计一次。这样各组无需调用组外通信资源，但确实增加了配置需求；它不是“免费共享”或缩短窗口后重新优化。若将中继架次也规定为不可复制，就要进一步合并服务区，属于另一种分区解释。

\subsection{区间峰值资源需求与精确枚举}
对任务组 $k$ 的某类同型资源 $h$，把占用区间记为 $\mathcal I_{kh}$，所需数量与库存缺口为
\begin{equation}
n_{kh}=\max_t\sum_{I\in\mathcal I_{kh}}\mathbf1\{t\in I\},\qquad
\Delta_h=\max\left\{0,\sum_kn_{kh}-I_h^{\rm stock}\right\}.\label{eq:partition}
\end{equation}
运输无人机占用至返航，运输电池延至充电完成；中继无人机延至返航并完成周转，中继能源组件延至返航并完成充电。与 Q2、Q3 相同的半开区间及时间容差避免把恰好衔接误判为额外需求。

峰值并发显然是资源数下界，因为同时占用的任务不能用同一资源。它也能达到：按开始时刻处理区间，将新区间交给已释放的同型资源；只有尚未释放的资源全部被占用时才增加一份，新资源数即此刻并发数。因此在允许组内同类资源重新编号、保持任务时刻不变的条件下，式\eqref{eq:partition}给出最少配置。这里保留任务安排，而非要求沿用 Q3 的实体资源编号。

对 $K=2$ 与 $K=3$，枚举无标号集合划分，候选数分别为 \PartitionsTwo{} 和 \PartitionsThree{}。规模很小，直接穷举能检查全部方案，无须再用启发式。对每个划分计算缺口、资源需求与组间工作量，并按
\begin{equation}
\operatorname{lexmin}\left(\sum_h\Delta_h,\ \sum_{k,h}n_{kh},\
\operatorname{CV}(w_k),\ \max_kN_k-\min_kN_k\right)\label{eq:q4objective}
\end{equation}
选择。$w_k$ 为运输占用时间加中继服务时间，$N_k$ 为运输架次数；$\operatorname{CV}$ 使用各组工作量的总体标准差除以均值。末项仅在前项相同时进一步平衡架次。资源数量直接相加是配置规模指标，官方未提供设备单价，故不解释为经济成本最小。

\subsection{峰值错位与中继独立配置的缺口分解}
记原共享计划的资源峰值为 $p_h$，分组独立总需求为 $g_h=\sum_kn_{kh}$，库存余量为 $u_h=I_h^{\rm stock}-p_h$。则
\begin{equation}
\Delta_h=\max\{0,(g_h-p_h)-u_h\}.
\end{equation}
运输资源的 $g_h-p_h$ 来自分组峰值错位，不能再利用另一组的空闲；中继资源的同一差额包含跨组复制完整中继任务的独立配置代价。先扣除共享计划未用库存，再得到真实缺口，所以两类增量直接相加通常不等于缺口。程序逐资源输出原峰值、独立需求、库存余量、增量和净缺口，并保存全枚举计数及Q3来源摘要。不同运输组批诱导的原子数不同，分法数量应随当前方案计算，不能为了匹配外部3280种而修改绑定关系。

\subsection{配置缺口、方案比较与小结}
\textbf{固定 Q3 方案独立分为两组和三组时，总资源缺口分别为 \GapTwo{} 和 \GapThree{}。}图\ref{fig:partition}把服务区归属与库存差异对应起来，表\ref{tab:resources}给出各类资源的精确需求。相同货箱仍被运输相同次数，但独立组内的峰值不能再与另一组的空闲相抵消。
\paperfig{fig10_partition}{取消跨组共享后，相同任务计划需要更多独立资源。黑色短竖线为库存，圆点与方点为两组、三组需求，红字为对应类型缺口；下图按不可拆连通分量排列服务区，用圆形、方形及标记内组号标识归属。按组复制中继完整服务安排的规则已计入需求，缺口不等于按价格加权的成本。}{partition}
\papertable{resource_gaps}{固定 Q3 计划的独立配置：组数增加后运输、电池与中继缺口分别核算}{resources}
两组方案保留更多组内复用机会；三组提升组织独立性，却增加当前库存不能覆盖的需求。现有某类资源即使有冗余，也不能抵消不同机型或不同用途的缺口，因此不宜只报告“总设备数充足”。验证器从 Q3 原始安排重新生成连通分量、占用区间与全部划分，并核对导出的组别和库存统计。

\subsection{均衡分区与最低资源偏好的区别}
上述资源优先方案的服务区数量分别为13/2和13/1/1，不能直接与参考论文的均衡分区比较。为作同口径偏好分析，在同一3280种划分上增加显式条件：两组各含6--9个服务区，三组各含4--6个；归一化运输工作量Gini分别不超过0.33和0.34。定义
\begin{equation}
G=\frac{\sum_{g,h}|W_g-W_h|}{2(K-1)\sum_gW_g},
\end{equation}
其中 $W_g$ 是准备至返航的运输占用总时长，中继忙时单列，不混入该公平性指标。另计算按货箱数加权的组内medoid离散度，服务区间阻抗取原DEM下B型机双向飞行时间平均值；中心仅用于衡量分区紧凑性，不改变出发基地O01。
\papertable{balanced_preferences}{同一Q3固定任务下，资源优先与均衡约束有不同配置代价}{balancedpreferences}
两组有118种、三组有414种划分满足均衡约束，最少缺口分别为\BalancedGapTwo{}和\BalancedGapThree{}；服务区分配为7/8和5/4/6，与参考论文缺口5/9达到同等汇总量级，但运输任务、中继复制及能耗不同。禁止跨组复制中继且每组至少两个服务区时，两组仅1种可行、三组无可行划分。偏好条件、可行计数、Gini与空间离散度的已枚举非支配集合均输出，资源优先方案仍保留供选择。均衡约束不是题设新增硬条件。

\textbf{本问小结。}Q4 完成了从联合调度到组织独立性的最后一步：在固定 Q3 时序、组内同型资源可重编号和中继完整任务可按组复制的解释下，精确求得分区配置。若允许重新组批、移动中继或调整时间，最少资源可能不同，本问不证明那个更大问题的全局最优。至此四问主线闭合，后续统一检查能力变化、紧余量及证据边界。
\FloatBarrier

\section{敏感性与固定方案压力边界}\label{sec:sensitivity}
四问得到的是官方确定性参数下的结果。为判断其可解释范围，须区分两种问题：改变参数后允许重新组批，最优目标怎样变化；保持现有计划不动，还能承受多少额外压力。前者反映重新决策后的适应性，后者只反映现有方案的余量，二者不能互相替代。本文只使用已计算的确定性场景，不构造未经观测的概率分布。

\subsection{储备变化与离散组批跳变}
储备提高直接收紧式\eqref{eq:feasible}，而不改变额定载荷。对 10\%--40\% 储备取 \ReserveGridCount{} 个等距点，重算服务区—机型安全载荷，共 \CapacityCaseCount{} 个能力场景；另在表\ref{tab:sensitivity}所列的 \BatchScenarioCount{} 个储备场景重新运行 Q1 精确组批。前者显示连续能力收缩，后者提供基准交叉核验。另枚举全部质量与体积可行子集，取各子集在官方机型中的最大可行储备，再用“最大化最小储备”的集合划分DP求每个精确架次数的储备上限；按所有临界点划分参数区间，并在各区间内部用原Q1求解器独立核对最少架次。因此图中的阶梯不是稀疏场景插值。
\paperfig{fig11_sensitivity}{提高储备使连续运载能力收缩，重新组批的离散架次数随之改变。上图细线为全部服务区，粗线追踪基准最受限服务区，不可达用轴下叉号单独标识；下图阶梯来自完整子集临界分析，空心点为原五场景独立重算，浅色斜纹区表示不可行；灰虚线分别为20\%基准和固定原方案储备上限。}{sensitivity}
\papertable{sensitivity}{储备增大对能力与组批的影响不同：表中每个场景均重新精确求解}{sensitivity}
基准及两个较低储备场景的组批目标保持一致，而更高储备下需要增加架次。原因不是算法突然失效，而是货箱不可连续切小：原先一个批次的固定质量越过能力边界时，必须换机型、拆成新的货箱子集或增加架次。因此，连续的物理参数变化会产生离散的组合决策跳变。

基准固定组批最低 SOC 为 \QOneReserveLimit\%，故不改变货箱、路线和机型时，最大共同储备就是该值。单凭这一固定方案值不能判断所有同架次数方案。本次额外的完整子集临界分析独立得到：最少架次首次在约 \ReserveTripBreak\% 后增加，该临界值在本算例恰与固定基准阈值重合；储备超过约 \ReserveReachabilityLimit\% 后，至少一个服务区不存在可行货箱划分。阶梯在临界点本身仍取较低架次数，图中端点和数值容差沿用统一能源可行判据。

\subsection{固定运输计划的能耗容忍度}
设所有运输架次的计算能耗统一乘以 $\lambda$，同时冻结路径、载荷与时间，则维持返航储备的最大系数为
\begin{equation}
\lambda_{\max}=\min_{p\in\mathcal P}
\frac{(1-\rho_{g_p})E_{g_p}}{E_p}.\label{eq:stress}
\end{equation}
每个架次分别给出一个上界，取最小值即可，无需新增随机仿真。Q2 与 Q3 当前运输计划的能耗上浮容忍度分别约为 \QTwoEnergyTolerance\% 与 \QThreeEnergyTolerance\%；Q3 最低 SOC 为 \MinSoc\%，仅略高于题设储备。因而“全部架次满足能源约束”只能证明基准可行，不能推断能耗估计有较大误差时仍可靠。

式\eqref{eq:stress}仅检查运输返航电量。它不代表整个联合方案可无条件承受同样增幅：能耗增加还可能延长电池充电，压缩后续任务可用时间；中继能源也需另行检查。本文没有把能耗放大后重新求解全部排程，因此不报告这种更强的恢复能力。

\subsection{通信余量与可解释范围}
保持链路几何、窗口和保障关系不变，若所有链路损耗统一增加同一常数，则在既有采样集合中可承受的增量上界等于最小可用余量。本次修复后的0.5 s回放最小值为 \MarginHalf{} dB，因而仅在固定采样集合和固定几何下允许同量级的共同额外损耗。正余量很小，不能据此推断实际衰落分布或可靠率。

新的地形遮挡可能改变链路可用模式，设备误差、天气和轨迹误差也不必等价于统一常数损耗。因此，此分析不能作为连续时间或新场景保证。若继续改进，应首先关注最紧通信区间和高能耗架次，再评估调整中继、分组和排程后的资源代价；本轮已修复已发现的空间链路反例，尚未进行上述误差情景下的全面重调度。
\FloatBarrier

\section{模型验证、评价与适用边界}\label{sec:validation}
本文的可信度来自可追溯的计算证据，而非算法名称。验证依次回答三个问题：计算是否针对同一套输入；输出是否真的满足规定约束；通过检查后还能声称多强的结论。四问的结果口径先汇总于表\ref{tab:summary}，其中累计作业与并行完工必须分别理解。
\papertable{summary}{四问使用同一物理口径；Q1 为累计作业，Q2 为运输完工，Q3 为联合完工，Q4 资源配置另见表\ref{tab:resources}}{summary}

\subsection{从数据版本到逐箱逐资源回放}
输入版本一致性核验防止将旧输出与新参数混用，也防止 Q4 从不同输入版本的 Q3 结果继续分区。论文诊断另记录正式输出摘要，生成宏、表和图时检查来源是否变化。这保证证据来自当前版本，但文件摘要本身不证明物理公式正确。

随后不直接相信结果表中的能耗、SOC 和交付时刻，而从官方箱单、访问顺序与机型重算：每箱恰好一次，卸货后更新载荷，空载返航，逐航段累计能量，交接结束才视为交付。资源核验分别检查无人机返航释放、电池充满释放、中继周转与组件充电；时限核验区分硬截止和软延误。Q2 保存的全部候选也逐一回放，Q4 则重新枚举固定计划的分区。

为检验验证器能否拒绝错误，使用输入篡改、输出能耗或时刻篡改、漏箱/重箱、嵌套占用区间与容差冲突等反例。物理边界涵盖零距离、零载荷、额定载荷、恰好达到储备、DEM 边界与角点、无效像元及双向链路覆盖内外。DEM 穿越另与独立矩形裁剪算法交叉核对，避免仅用同一几何函数产生测试预期。

这些证据有明确分工：回放检查约束和表格真实性，反例检查错误是否会被发现，独立几何参照与量纲核对补充共享内核的盲区。能耗回放仍调用官方物理引擎，所以不能排除所有共享公式错误，更不能代替实际飞行校准。具体测试命令和完整记录随源码保存，正文只保留支持模型可信范围的证据。

\subsection{加密通信回放与最紧约束}
\paperfig{fig12_validation}{基准可行并不等于余量充足：最低返航 SOC 为 \MinSoc\%，最小硬截止余量为 \HardSlack{} s，最小通信余量为 \MarginHalf{} dB。三类证据分别对应能源、及时性和通信风险。}{validation}
\papertable{sampling}{修复后的同一联合计划在1 s与0.5 s回放中均无未覆盖状态}{sampling}
图\ref{fig:validation}和表\ref{tab:sampling}给出最紧约束。\HardBoxes{} 个硬截止货箱均未违反截止，其中最紧余量为\HardSlack{} s。通信 1 s 设置检查 \StatesOne{} 条记录，0.5 s 设置检查 \StatesHalf{} 条，两者均无未覆盖状态，最小余量分别为 \MarginOne{} 与 \MarginHalf{} dB。另以未参与搜索的0.25 s网格回放 \StatesQuarter{} 条状态，未覆盖数为 \ViolationsQuarter{}；该检查通过仍不代表连续时间证明。

两种网格都按各相位长度分别等分，并不严格嵌套；状态记录还包含相位边界重复点，所以既不能称为独立统计样本，也不能据两点画出一般性的收敛规律。遮挡判定可能在采样点之间突变。连续通信安全仍需要定位遮挡变化事件，或为整段时间建立可验证的余量下界；当前证据尚未完成这一步。因此，Q3 面对题设连续通信要求仍存在明确证明缺口。

\subsection{四类结论的精确性范围}
四问所称“解”的含义不同，表\ref{tab:scope}集中列出其边界。特别是，精确求解一个受限问题，不等于精确求解包含所有离散和连续选择的原问题。
\begin{table}[htbp]\centering\caption{算法证据与结论强度一一对应}\label{tab:scope}
\begin{tabularx}{\linewidth}{lXX}\toprule
问题 & 已成立的结论 & 不由当前结果证明的结论\\\midrule
Q1 & 单点往返、当前物理规则与词典序目标下，集合划分 DP 精确 & 加入实体资源、时限及多点访问后的调度最优\\
Q2 & 架次列与区间着色生成全准时候选；有限目录和网格上提供修正对偶下界 & 全局最优、完整 Pareto 前沿、所选架次数唯一最优\\
Q3 & 外层离散搜索与条件时序LP；1 s、0.5 s及独立0.25 s回放通过 & 路线、货箱、中继位置/高度的联合全局最优及连续通信安全\\
Q4 & 固定 Q3 时序及中继复制解释下，分区枚举精确 & 允许全面重规划时的全局最少资源或经济成本最优\\\bottomrule
\end{tabularx}\end{table}
\FloatBarrier

\subsection{模型优势、局限与改进顺序}
公共物理与上层决策分离，使每个结果均可沿货箱、航段、资源占用和通信状态追溯。Q1与固定Q3后的Q4可精确求解；Q2和Q3采用受限列模型、明确热启动与固定种子邻域搜索，没有原问题全局最优或统计显著性证据。官方等效能耗未经实飞辨识，也不支持天气外推。

已发现的通信反例经中继修复与重新搜索消除。后续应优先定位采样间的遮挡变化；再评估能耗误差、充电变化与最紧截止余量下的可执行性，最后扩大分组、中继与资源顺序搜索。现有候选和反例均可复算，不能替代实飞验证。
\FloatBarrier

\section{结论}
本文以统一的地形、动态载荷、能量与资源占用模型连接四问。Q1用精确集合划分得到\QOneTrips{}架次；Q2多邻域搜索按延误优先选出\QTwoTrips{}架次；Q3外层搜索与条件LP形成\QThreeTrips{}运输、\RelayTrips{}中继的正式候选；Q4固定其任务时序，精确得到两组、三组分别缺少\GapTwo{}、\GapThree{}份资源。延误改善伴随运输架次和能耗增加，不是所有指标同时改善，外层搜索也没有全局最优证书。

独立加密验证限定了结论：Q3在1 s采样下通过，0.5 s却发现\ViolationsHalf{}个未覆盖状态，最低余量\MarginHalf{} dB；最小硬截止余量为零。当前联合候选尚需修复通信缺口，Q4仅为其固定输出下的条件配置分析。后续应优先补足沿程通信验证与覆盖修复，再扩大离散搜索和开展真实能耗校准。

\clearpage
\begin{thebibliography}{9}
\bibitem{problem} 中国研究生数学建模竞赛组委会. 2026 年中国研究生数学建模竞赛 D 题：山区洪涝灾害下无人机运输与通信协同优化及配套数据[Z]. 2026.
\bibitem{karney} Karney C F F. Algorithms for geodesics[J]. Journal of Geodesy, 87(1): 43--55, 2013. DOI: \url{10.1007/s00190-012-0578-z}.
\bibitem{dorling} Dorling K, Heinrichs J, Messier G G, Magierowski S. Vehicle Routing Problems for Drone Delivery[J]. IEEE Transactions on Systems, Man, and Cybernetics: Systems, 47(1): 70--85, 2017. DOI: \url{10.1109/TSMC.2016.2582745}.
\bibitem{zhang} Zhang J, Campbell J F, Sweeney D C, Hupman A C. Energy consumption models for delivery drones: A comparison and assessment[J]. Transportation Research Part D: Transport and Environment, 90: 102668, 2021. DOI: \url{10.1016/j.trd.2020.102668}.
\bibitem{zeng} Zeng Y, Zhang R, Lim T J. Wireless communications with unmanned aerial vehicles: Opportunities and challenges[J]. IEEE Communications Magazine, 54(5): 36--42, 2016. DOI: \url{10.1109/MCOM.2016.7470933}.
\bibitem{itu} International Telecommunication Union. Recommendation ITU-R P.525-5: Calculation of free-space attenuation[S/OL]. 2024[2026-09-23]. \url{https://www.itu.int/rec/R-REC-P.525-5-202411-I/en}.
\bibitem{highs} Huangfu Q, Hall J A J. Parallelizing the dual revised simplex method[J]. Mathematical Programming Computation, 10(1): 119--142, 2018. DOI: \url{10.1007/s12532-017-0130-5}.
\end{thebibliography}

\clearpage
\appendix
\section{完整结果与复现说明}
\subsection{安全载荷与单点组批}
本附录表格保留论文评阅需要的完整服务区与架次结果。货箱逐项分配、逐箱交付、通信保障区间和任务分区的完整可机读记录见随附官方格式工作簿与 CSV。数值显示精度不改变源结果的完整精度。
\longresulttable{payloads}{15 个服务区的最大安全载荷；基准返航储备为 20\%}{payloads}
\longresulttable{q1_batches}{Q1 全部单点架次；表中前缀 Q1 省略}{q1batches}
\longresulttable{frontier}{Q2 全部候选的数值指标}{frontier}
\subsection{运输与中继完整架次表}
电池简写 A-01 表示 A-BAT-01，其余同理。所有时刻从任务起点计，以 s 为单位。实际访问顺序和货箱列表保留在工作簿，避免把过宽字符串压缩为不可读的小字号。
\longresulttable{q2_schedule}{Q2 全部运输架次；架次前缀 Q2 省略}{q2schedule}
\longresulttable{q3_schedule}{Q3 全部运输架次；架次前缀 Q3 省略}{q3schedule}
\longresulttable{relay_schedule}{中继飞行、服务窗口与能耗}{relayschedule}
\longresulttable{relay_positions}{中继悬停位置与返航 SOC}{relayposition}
\subsection{计算入口与证据文件}
将官方附件原样解压至程序的数据目录后，依次运行以下命令可生成正式结果、诊断和图表。所有路径在项目根目录下解析，Windows 与 Linux 使用同一 Python 入口；论文需要含中文支持的 XeLaTeX。
\begin{verbatim}
python run_all.py
python -m pytest -q
python scripts/analyze_paper.py
python scripts/setup_paper_fonts.py
python scripts/build_paper.py
\end{verbatim}
正式结果验证文件记录官方输入的 SHA256 摘要。图形契约记录每图的核心主张、数据路径、脚本、单位与不确定性口径；所有主图同时提供 PDF、SVG 和 300 dpi PNG。新增采样诊断不覆盖正式结果，且每次纸面数字更新都应重新生成图、表和工作簿。

\section{人工智能辅助与第三方来源说明}
本工作使用 ChatGPT/Codex（OpenAI，2026-09-23）辅助代码审计、图形设计和文字草拟。提示要求以官方规则为准、不编造数据与文献；输出经数据回放、边界测试、加密通信检查和逐页审查。本次精确模型版本与公开发布日期在执行记录中不可得，未沿用无法核验的历史声明；参赛队应从实际使用界面补全官方要求的信息，并独立核验和理解全部内容。AI 建议不作为公式或实验数据来源。

科研图形组织参考 TAO-QKV 的 Icarus-Figures（\url{https://github.com/TAO-QKV/Icarus-Figures}，访问日期 2026-09-23）：采用一图一主张、突出主面板、展示约束阈值、灰度冗余编码和三格式导出。全部无人机场景图由本项目真实数据重新绘制，没有粘贴外部示例图片，也未使用生成式图像。

\end{document}
